\originalchapter{感应、自感与互感}\label{ch:induction}
\originalsection{法拉第定律与楞次方向}
静止线圈中的磁铁移动时, 即使没有电池也会出现电流. 与现象相联系的量是磁通$\Phi_B=\int_S\bb\cdot\dd\vct A$, 单位韦伯(Wb). 固定回路的感应电动势满足
\[
\mathcal E=\oint_C\ee\cdot\dd\vct\ell=-\frac{\dd\Phi_B}{\dd t}.
\]
环路与法向依右手规则关联. 负号表示感生电流倾向阻碍原磁通变化, 这叫楞次定律. 阻碍的是“变化”, 不总是与原磁场反向: 原向外磁通减少时, 感生场向外以补偿减少.

移动回路还受$q\vct v_{\rm wire}\times\bb$驱动, 总电动势应写$\oint(\ee+\vct v_{\rm wire}\times\bb)\cdot\dd\vct\ell=-\dd\Phi_B/\dd t$. 此处右边是随回路运动的通量变化率. 固定回路中的变化磁场产生非保守电场, 因而不能在整条回路上都用单值静电势差替代感应电动势.

\begin{example}
一根长$\ell$导棒沿平行导轨以速度$v$滑动, 回路电阻$R$, 均匀磁场$B$垂直轨道平面. 求维持匀速所需功率.
\end{example}
\begin{solution}
回路面积变化率为$\ell v$, 所以$|\mathcal E|=B\ell v$, $I=B\ell v/R$. 导棒磁力大小$IB\ell=B^2\ell^2v/R$, 方向阻碍运动. 外力须同样大小而顺运动, 功率$P=Fv=B^2\ell^2v^2/R=I^2R$. 若不施外力, 导棒动能会转为焦耳热. 忽略自感时该关系即时成立; 自感明显时还要计磁场储能变化.
\end{solution}

绕固定轴旋转的$N$匝线圈, 均匀场中磁通链为$NBA\cos\omega t$, 感应电动势$NBA\omega\sin\omega t$. 这给出发电机的正弦驱动. 发电不是能量创造, 维持转动的机械功经电磁过程送到负载. 实体金属内部的涡流也遵循该规律, 常导致磁阻尼与热损耗.

\originalsection{磁通链与电感}
同一线圈的各匝磁通相加为磁通链$\Lambda=\sum_i\Phi_i$. 在线性、几何固定模型中$\Lambda=LI$, $L$称自感, 单位亨利(H), 于是$\mathcal E_L=-L\dot I$. 电流增大时感应作用反对增大, 电流减小时反对减小. 电感元件两端按被动符号约定的电压降是$V_L=L\dot I$, 不应把这个电压降与回路内自感电动势的负号混淆.

长螺线管内部$B=\mu_0nI$, 每匝通量为$BA$, 匝数$N=n\ell$, 故$L=\mu_0n^2A\ell$. 电源建立电流时对抗自感的功率$P_L=IL\dot I$, 积分得到$U_B=LI^2/2$. 代入螺线管场与体积给出真空磁能密度$u_B=B^2/(2\mu_0)$. 非线性磁芯需改为实际积分关系, 不能保留常数$L$而任意提高电流.

两线圈之间的磁通链可写$\Lambda_2=M_{21}I_1$, 所以线圈2的感应电动势$\mathcal E_2=-M_{21}\dot I_1$. 在通常互易线性介质中$M_{12}=M_{21}=M$, 正负随线圈绕向和端点约定改变. 理想变压器两线圈链接同一磁通, $|V_2/V_1|=N_2/N_1$. 忽略损耗时$|V_1I_1|=|V_2I_2|$, 电流比反向; 恒定直流在暂态结束后不产生持续互感电动势.

同轴电缆内外导体间$B=\mu_0I/(2\pi s)$, 单位长度场能为$\int_a^b[B^2/(2\mu_0)]2\pi s\dd s=\mu_0I^2\ln(b/a)/(4\pi)$. 与$L'I^2/2$比较, 得外部电感每单位长$L'=\mu_0\ln(b/a)/(2\pi)$. 内导体内若仍有磁场需另加内部电感, 不能把外部电感当成任意电缆总电感.


金属盘进入或离开局部磁场区时, 不同盘内回路的磁通变化, 产生环流. 环流在外场中受力阻碍盘转动, 外力维持转动所供机械能成为焦耳热; 不接外线的实体金属也能构成这些闭合回路. 相同输电功率$P$, 电压升高使电流$I=P/V$降低, 输电线热损$I^2R_{\rm line}=P^2R_{\rm line}/V^2$下降. 理想变压器由匝数比改变交流电压和电流, 并不提高输入总功率; 实际变压器还存在铜损、铁损与绝缘约束.
